% ---------------------------------------------------------------------
%  Chapter 17 : The compute_cumulants Orchestrator
% ---------------------------------------------------------------------
\chapter{The \texttt{compute\_cumulants} Orchestrator}\label{ch:compute-orchestration}

Every chapter before this one has opened up a single subsystem in isolation: the
field-theory core that expands the action, the propagator builder, the
mean-field solver, the diagram enumerator, the symmetry-factor machinery, Phase~J
in the time domain, the spatial integrator. This chapter is where they are
\emph{wired together}. There is exactly one function in the codebase that calls
all of them in order, threads the data from each into the next, and hands you
back a finished correlation function: \code{compute\_cumulants}, defined in
\file{pipeline/compute.py:108}.

If the rest of the manual is the anatomy of the engine, this chapter is its
nervous system. \code{compute\_cumulants} is the function the front-desk module
\code{daedalus.dd.run} ultimately calls (you saw that call in
Chapter~\ref{ch:quickstart}, at \file{notebooks/daedalus.py:586}); it is the
function the simulation-comparison harnesses call; it is the function the example
notebooks narrate when they print their seven-phase trace. Understanding it is
understanding how a model declaration becomes a number.

The crucial thing to hold in mind from the first page is that
\code{compute\_cumulants} \emph{does almost no arithmetic of its own.} It is
\term{orchestration}, not computation. Its job is to sequence seven phases,
carry a handful of data dictionaries between them, decide whether the model is
temporal or spatial and branch accordingly, assemble the per-loop-order pieces
into a master correlator, and stamp wall-time onto each phase so you can see
where the seconds went. The heavy lifting all happens in the sibling modules
this chapter's code merely calls. So this chapter documents two things in
parallel: the \emph{glue} (what \code{compute\_cumulants} itself does, line by
line) and the \emph{contracts} (what each phase promises to return, so the next
phase can consume it). Where we name a function from another module we cite it
with its file and line so you can jump straight to its own chapter.

\section{Where this function sits in the pipeline}\label{sec:co-where}

Recall the spine from Chapter~\ref{ch:quickstart}. The user writes a
\term{theory file}, \code{dd.load\_theory} turns it into a \code{model} dict, and
\code{dd.run} packages the run choices and calls the engine. The engine is this
function. End to end:

\begin{lstlisting}[style=console]
  model dict  (from a TheoryBuilder / hand-written models/*.py)
       |
       |  + fundamental (numeric params)  + k, max_ell, external_fields
       v
  +-----------------------------------------------------------------+
  |          compute_cumulants   (pipeline/compute.py:108)          |
  |                                                                 |
  |  [1] expand      FieldTheory.expand()      -> bigraded action   |
  |  [2] propagator  build_propagator()        -> K_ft, G_ft, ...   |
  |  [3] mean_field  solve_mean_field[_dae]()  -> num_params, saddle|
  |  [3.5] spatial?  -> momentum integrator (early return)          |
  |  [4] poles       compute_poles_and_residues() -> pole_vals,Cmat |
  |  [5] diagrams    enumerate_unique_diagrams() -> typed diagrams  |
  |  [6] classify    classify_coefficient_factors() -> prefactors   |
  |  [7] phase_j     compute_correction_td() per ell -> C(tau)      |
  +-----------------------------------------------------------------+
       |
       v
  result dict  -> notebooks (plotting), pipeline.save (npz/csv),
                  sim-comparison harnesses
\end{lstlisting}

\textbf{What feeds it.} The \code{model} dict (the subject of the
\emph{theory-files} chapters), built either by a \code{TheoryBuilder} or written
by hand under \file{models/}. Alongside it: the numerical parameter point
\code{fundamental} (a dictionary mapping each coupling's name to its value), the
cumulant order \code{k} (the number of external legs), the maximum loop order
\code{max\_ell}, and \code{external\_fields} (which field sits on each external
leg).

\textbf{What consumes its output.} The returned \code{result} dict. Its
\code{total\_C} callable and its \code{C\_tau} / \code{C\_tau\_by\_ell} arrays
are what the notebooks plot and what the simulation harnesses compare against;
\code{pipeline.save.save\_npz} and \code{save\_csv} serialize it; the \code{mf}
and \code{params} accessors expose the saddle values and parameters by their
natural names. We dissect that dict in Section~\ref{sec:co-result}.

\begin{note}
The seven labels \code{[1/7]} through \code{[7/7]} that the engine prints in a
verbose run --- you saw them in Chapter~\ref{ch:quickstart} --- are emitted by
this file, one \code{print} per phase. The whole rest of this manual is, in a
sense, a deep dive into one or more of those seven labels. This chapter is the
table of contents that ties them together.
\end{note}

\section{The \protect\msrjd{} expansion in one page}\label{sec:co-msrjd}

To read the orchestrator you need only the skeleton of the physics; the
field-theory chapters fill it in. A stochastic dynamics for a field $\varphi$ ---
a firing rate, a density, a height field --- is recast as a path integral over
$\varphi$ \emph{and} a conjugate \term{response field} $\tilde\varphi$ (the
``tilde'' field). The generating functional is
\begin{equation}
  Z[j,\tilde\jmath]
  \;=\;
  \int \mathcal{D}\varphi\,\mathcal{D}\tilde\varphi \;
  \exp\!\Big(-S[\varphi,\tilde\varphi]
    + \!\int (\tilde\jmath\,\varphi + j\,\tilde\varphi)\Big).
  \label{eq:co-Z}
\end{equation}
This is the \term{Martin--Siggia--Rose / Janssen--De Dominicis} (\msrjd)
construction. The action $S$ splits into three pieces:

\begin{itemize}
  \item a \term{free (Gaussian, bilinear) part} $S_0$, quadratic in the fields;
        its inverse is the bare propagator;
  \item an \term{interaction part} $S_{\mathrm{int}}$, the cubic and higher
        vertices;
  \item a \term{mean-field / saddle part}, the linear and constant terms, which
        vanish exactly when the expansion is taken about the correct saddle
        point.
\end{itemize}

Connected correlation functions --- \term{cumulants} --- are obtained by
functionally differentiating $\ln Z$. The key fact for everything that follows is
that the connected $k$-point cumulant is precisely the sum of all
\emph{connected} Feynman diagrams with $k$ external legs. \code{compute\_cumulants}
automates that sum: it expands the action, builds the propagator, solves for the
saddle, enumerates every diagram up to \code{max\_ell} loops, computes each
diagram's symmetry factor, and integrates each one.

\begin{defn}
\textbf{Bigrade.} Each monomial in the expanded action carries a \term{bigrade}
$(n_t, n_p)$ = (number of response-field $\tilde\varphi$ factors, number of
physical-field $\varphi$ factors). The bigrade is how the expanded action is
sorted into sectors:
\begin{itemize}
  \item \textbf{Mean-field sector} --- bigrades $(0,0)$, $(1,0)$, $(0,1)$. These
        linear/constant terms vanish at the saddle.
  \item \textbf{Free action} --- bigrade $(1,1)$, the bilinear kernel. Its
        time-domain matrix is \code{K\_ker}; its Fourier image is \code{K\_ft};
        its inverse \code{G\_ft = K\_ft$^{-1}$} is the bare propagator.
  \item \textbf{Interaction sectors} --- total degree $\ge 2$ excluding $(1,1)$.
        These are the vertices of the diagrammatic expansion.
\end{itemize}
\end{defn}

The helper \code{\_print\_action\_sectors} (\file{pipeline/compute.py:50}) prints
exactly these three groups after the expand step, when \code{verbose=True}. It is
the ``here is your action'' diagnostic and has no effect on the computation; we
return to it in Section~\ref{sec:co-helpers}.

\section{The seven phases at a glance}\label{sec:co-glance}

The body of \code{compute\_cumulants} is, at heart, a straight-line sequence of
seven numbered phases with a single spatial-versus-temporal fork wedged between
phases [3] and [4]. Here is the whole tour in one table; the sections that follow
open each row.

\begin{center}
\small
\begin{longtable}{@{}p{0.10\textwidth}p{0.32\textwidth}p{0.50\textwidth}@{}}
\toprule
\textbf{Phase} & \textbf{Calls} & \textbf{Produces} \\
\midrule
\endhead
{}[1] expand     & \code{FieldTheory.expand()} & the bigraded action \code{ft} (free, interaction, MF sectors) \\
{}[2] propagator & \code{build\_propagator()}  & \code{prop} dict (\code{K\_ft}, \code{G\_ft}, \code{D\_delta}, $\omega$, \code{nf}) \\
{}[3] mean field & \code{solve\_mean\_field[\_dae]()} & \code{num\_params}: the $\{$\code{SR}$\to$float$\}$ substitution dict \\
{}[3.5] spatial? & spatial integrator & \emph{early return} for spatial models \\
{}[4] poles      & \code{compute\_poles\_and\_residues()} & fills \code{prop['pole\_vals']}, \code{prop['C\_mats']} \\
{}[5] diagrams   & \code{enumerate\_unique\_diagrams()} & \code{unique\_by\_ell}: $\{\ell:$ typed diagrams$\}$ \\
{}[6] classify   & \code{classify\_coefficient\_factors()} & \code{diagram\_records}: prefactor + $\Sym(\Gamma)$ per diagram \\
{}[7] phase J    & \code{compute\_correction\_td()} per $\ell$ & \code{C\_tau\_by\_ell}, \code{total\_C\_by\_ell} \\
\bottomrule
\end{longtable}
\end{center}

The data threads through linearly: \code{ft} from [1] feeds [2], [3], and [5];
the \code{prop} dict from [2] is filled in place by [4] and consumed by [7];
\code{num\_params} from [3] is the universal substitution dictionary every
numerical phase needs. We trace that flow explicitly in
Section~\ref{sec:co-dataflow}.

\section{Argument validation and the Taylor budget}\label{sec:co-args}

Before any phase runs, the function validates its arguments and picks a few
derived quantities. The full signature
(\file{pipeline/compute.py:108}) is long because the function is the single
public surface for a great many run options:

\begin{lstlisting}
def compute_cumulants(
    model: dict,
    k: int,
    max_ell: int = 0,
    fundamental: dict = None,
    external_fields: list[tuple[str, int]] = None,
    *,
    tau_max: float = 50.0,
    tau_step: float = 0.5,
    tau_grid=None,                   # explicit tau grid; overrides tau_max/tau_step
    spatial_grid=None,
    taylor_order: int = None,
    origin_leaf_idx: int = 0,
    output_npz: str = None,
    output_csv: str = None,
    use_cache: bool = True,
    parallel: bool = True,
    spatial_parallel: bool = True,
    spatial_n_q: int = 64,
    spatial_points=None,
    n_workers: int = None,
    use_grouped_phase_j: bool = False,
    fixed_point_index: int = 0,
    mf_dae_n_starts: int = 64,
    mf_dae_seed_box: dict | None = None,
    verbose: bool = True,
) -> dict[str, Any]:
\end{lstlisting}

The load-bearing arguments are the first five: \code{model}, \code{k},
\code{max\_ell}, \code{fundamental}, and \code{external\_fields}. The keyword-only
block (\code{*}) after them tunes the run --- grid extents, caching, parallelism,
the spatial APIs, the DAE multi-root solver --- and all have defaults.

\textbf{Validation} (\file{pipeline/compute.py:253}). The function defaults a
missing \code{fundamental} to \code{\{\}}, then enforces the one structural
invariant it can check up front: \code{external\_fields} must be present and must
have exactly \code{k} entries.

\begin{lstlisting}
if fundamental is None:
    fundamental = {}
if external_fields is None:
    raise ValueError('external_fields is required')
if len(external_fields) != k:
    raise ValueError(
        f'external_fields has {len(external_fields)} entries but k={k}'
    )
\end{lstlisting}

\subsection{Why \texorpdfstring{$k + 2\ell$}{k + 2 ell} vertices, and the floor of 2}

The nonlinear functions in an action (a sigmoid firing rate, a quartic
restoring force) are Taylor-expanded to a finite order. How many terms do we
need? A connected diagram with $k$ external legs at $\ell$ loops needs vertices
of degree at most $k + 2\ell$: every loop you add introduces two extra legs that
have to terminate somewhere, and the worst case is that they all terminate on a
single vertex. So the smallest sufficient Taylor order is

\begin{equation}
  \texttt{taylor\_order} \;=\; \max(k + 2\,\texttt{max\_ell},\ 2),
  \qquad\text{(\file{pipeline/compute.py:279})}
  \label{eq:co-taylor}
\end{equation}

implemented as:

\begin{lstlisting}
if taylor_order is None:
    taylor_order = max(k + 2 * max_ell, 2)
\end{lstlisting}

The floor of $2$ is structural, not a heuristic. Order $2$ is the minimum that
still produces the bilinear $(1,1)$ propagator kernel and the
$(0,0)/(1,0)/(0,1)$ mean-field saddle sectors --- and downstream code reads those
\emph{unconditionally}, even in the degenerate $k=2,\ \texttt{max\_ell}=0$ case
where the tree-level pair correlator is just the bare propagator with no
interaction vertices at all.

\begin{note}
The docstring (\file{pipeline/compute.py:179}) records that this floor was
\emph{historically 4}, chosen for a now-obsolete cache layout that kept all
1-loop 2-cumulant runs in one directory. Lowering it from $4$ to $2$ saved about
$90$ minutes on heavy Bernoulli theories at $k=2,\ \texttt{max\_ell}=0$, the one
case where the old floor exceeded the mathematical minimum. If you are probing
higher-order vertices for cache-invalidation testing you must now pass an
explicit \code{taylor\_order} override.
\end{note}

\subsection{Natural names to internal names}

Users pass external fields by their natural names, e.g. \code{('n', 1)} for the
first population's physical field. The action and propagator-typing code, by
contrast, work with internal \emph{fluctuation} names, e.g. \code{('dn', 1)}. The
translation happens at \file{pipeline/compute.py:288}:

\begin{lstlisting}
naming_convention    = model.get('naming_convention')
external_fields_user = list(external_fields)
external_fields      = normalize_external_fields(
    external_fields, naming_convention=naming_convention)
\end{lstlisting}

\code{normalize\_external\_fields} (in \file{pipeline/access.py}) reads the
model's declared \code{naming\_convention} --- written by the
\code{TheoryBuilder} --- to map \code{('n',1)} to \code{('dn',1)}, and falls back
to a classic n/v/m map when the model declares none. Notice that the function
keeps the \emph{user's} original form in \code{external\_fields\_user}; both forms
end up in the result's \code{config} block (the internal form under
\code{external\_fields}, the user form under \code{external\_fields\_in}) so a
caller can always recover what it asked for.

\subsection{The per-phase wall-time tracker}

The last thing set up before phase [1] is the timing apparatus
(\file{pipeline/compute.py:296}):

\begin{lstlisting}
phase_walls: dict[str, float] | None = {} if verbose else None

def _phase_time(label: str, t0: float) -> None:
    if phase_walls is None:
        return
    dt = time.perf_counter() - t0
    phase_walls[label] = dt
    print(f'      [{label}] done in {dt:.2f}s')
\end{lstlisting}

\code{phase\_walls} is a dict mapping a phase label to its elapsed seconds, or
\code{None} when \code{verbose=False} --- so the bookkeeping is opt-in and costs
nothing in a quiet run. \code{time.perf\_counter()} is Python's highest-resolution
monotonic clock (it does not jump if the system clock is adjusted), which is why
it is the right tool for measuring elapsed wall-time rather than \code{time.time()}.
The closure \code{\_phase\_time(label, t0)} is called at the end of each phase
with that phase's start time; the per-phase summary at the very end of the
function (Section~\ref{sec:co-perf}) reads this same dict.

\section{Phase [1] --- expand, with the on-disk cache}\label{sec:co-expand}

\textbf{Code:} \file{pipeline/compute.py:305}--\code{:351}.

The first phase constructs the \code{FieldTheory} object and expands the action.
Expansion --- the multivariate Taylor pass that turns the symbolic action into
the bigraded sector dictionary --- is one of the more expensive symbolic
operations in the pipeline, so it is cache-aware:

\begin{lstlisting}
ft = FieldTheory(model, taylor_order=taylor_order)

from pipeline import _expand_cache as _ec
cache_hit = False
if use_cache:
    cached_order = _ec.find_best_cached_order(model, taylor_order)
    if cached_order is not None:
        _ec.prepare_for_load(ft)
        cache_hit = _ec.load_expand(
            model, ft,
            target_order=taylor_order,
            cached_order=cached_order,
            verbose=verbose,
        )
if not cache_hit:
    ft.expand()
    if use_cache:
        try:
            _ec.save_expand(model, ft, verbose=verbose)
        except Exception as e:
            ...
\end{lstlisting}

The logic is worth spelling out because it is more clever than a plain
hit-or-miss cache. \code{find\_best\_cached\_order} looks for \emph{any} cached
expansion at an order \emph{$\ge$} the one requested --- because an order-6
expansion already contains every term of an order-4 expansion. If it finds one,
\code{prepare\_for\_load} and \code{load\_expand} downgrade-filter the cached
sectors to the requested total degree and load them. Only on a genuine miss does
it pay for \code{ft.expand()} and then \code{save\_expand} the result. The
\code{try/except} around the save is deliberate: a failed cache write is logged
and ignored, never fatal --- you would rather finish the run uncached than crash
on a read-only disk.

After expansion the code runs a self-check and extracts the vertex/source types:

\begin{lstlisting}
sanity_ok = ft.sanity_check(verbose=verbose)
if not sanity_ok:
    raise RuntimeError(
        'FieldTheory.sanity_check() failed -- see printout above.')

vtypes = extract_vertex_types(ft)
stypes = extract_source_types(ft)
n_noise = sum(1 for s in stypes if isinstance(s, NoiseSourceType))
\end{lstlisting}

\code{vtypes} are the interaction vertices the enumerator will draw from;
\code{stypes} are the source (noise) vertices, and \code{n\_noise} counts how
many are genuine \code{NoiseSourceType}s. These three are produced here and
carried all the way to phase [5]. When verbose, the phase ends by calling
\code{\_print\_action\_sectors(ft)} (the diagnostic of
Section~\ref{sec:co-msrjd}) and records its wall-time with
\code{\_phase\_time('expand', \_t\_phase)}.

\section{Phase [2] --- the propagator}\label{sec:co-prop}

\textbf{Code:} \file{pipeline/compute.py:353}--\code{:358}.

The propagator is the inverse of the free action's $(1,1)$ kernel. This phase is
a single call:

\begin{lstlisting}
prop = build_propagator(ft, model, use_cache=use_cache, verbose=verbose)
\end{lstlisting}

\code{build\_propagator} (\file{pipeline/\_propagator.py:524}) builds (or loads
from \file{saved\_theories/<tag>/propagator.sobj}) the symbolic propagator
dictionary \code{prop}. Its keys are \code{K\_ker}, \code{K\_ft}, \code{G\_ft},
\code{adj\_ft}, \code{D\_omega}, \code{D\_delta}, \code{t\_var}, \code{omega},
\code{nf}, and \code{ring\_gen\_names}, plus \code{pole\_vals} and \code{C\_mats}
placeholders that phase [4] will fill. The complete anatomy of that dict is in
the propagator chapter and Section~\ref{sec:co-datastructures}; the orchestration
view is just: \emph{this is where the bare two-point function comes from.}

\subsection{The math the propagator chapter spells out}

The free action's $(1,1)$ sector is a matrix \code{K\_ker}$(t)$ of time-domain
kernels (containing Dirac $\delta(t)$, its derivative $\delta'(t)$, and any model
kernel symbols). Fourier-transforming gives \code{K\_ft}$(\omega)$, a matrix of
rational functions of $\omega$. The bare retarded propagator is its matrix
inverse:
\begin{equation}
  \text{\code{G\_ft}}(\omega) \;=\; \text{\code{K\_ft}}(\omega)^{-1}
  \qquad\text{(rows = physical, cols = response).}
  \label{eq:co-Gft}
\end{equation}
Its poles in the upper-half $\omega$-plane are the retarded relaxation modes; the
residue matrices \code{C\_mats} at those poles give the time-domain propagator as
a sum of decaying exponentials,
\begin{equation}
  G(t)
  \;=\;
  \sum_k C_k\, \ee^{-\ii \omega_k t}\,\Theta(t)
  \;+\;
  \texttt{D\_delta}\,\delta(t),
  \label{eq:co-Gt}
\end{equation}
where \code{D\_delta}$= \lim_{\omega\to\infty}$\code{G\_ft}$(\omega)$ is the
instantaneous (white) part. Phase [2] produces the symbolic \code{G\_ft} and
\code{D\_delta}; phase [4] produces the numbers $\omega_k$ (in \code{pole\_vals})
and $C_k$ (in \code{C\_mats}). For ``rich'' kernels --- large matrices, many
symbols --- \code{build\_propagator} deliberately skips the symbolic inverse and
leaves \code{G\_ft = None}, deferring the whole job to the numerical pole-finder
in phase [4].

\section{Phase [3] --- the mean-field saddle}\label{sec:co-mf}

\textbf{Code:} \file{pipeline/compute.py:360}--\code{:380}.

The loop expansion is taken about the deterministic mean-field steady state. This
phase finds it and --- more importantly for the orchestrator --- assembles the
substitution dict that makes every later phase numeric:

\begin{lstlisting}
if model.get('equations'):
    mf = solve_mean_field_dae_compat(
        ft, model, fundamental,
        fixed_point_index=fixed_point_index,
        n_starts=mf_dae_n_starts,
        seed_box=mf_dae_seed_box,
        verbose=verbose,
    )
else:
    mf = solve_mean_field(ft, model, fundamental, verbose=verbose)
num_params = mf['num_params']
\end{lstlisting}

There are two solvers, chosen by whether the model declares \code{equations}:

\begin{itemize}
  \item \code{solve\_mean\_field} (\file{pipeline/\_mean\_field.py:14}) is the
        legacy / single-population path. It assembles the parameter substitutions
        from \code{fundamental}, reads the symbolic background $v^{*}$ from the
        model's \code{mf\_bg\_conditions}, and solves the self-consistency
        $n^{*}_i = \langle\,\text{RHS}\,\rangle(n^{*}_j)$ numerically.
  \item \code{solve\_mean\_field\_dae\_compat}
        (\file{pipeline/\_mean\_field\_dae.py}) is the DAE-based multi-root
        solver, used when the theory declares its dynamics via
        \code{.equation(lhs=..., rhs=..., population=...)}. It layers a
        multi-start global root search on top, returning \emph{every} fixed point
        with a stability annotation, and \code{fixed\_point\_index} selects among
        the stable ones.
\end{itemize}

\begin{defn}
\textbf{\code{num\_params}.} \code{num\_params} is a dictionary
$\{$\code{SR.var} $\to$ \code{float}$\}$. It
maps every symbolic parameter, every solved saddle value (\code{nstar\_i},
\code{vstar\_i}, \code{mstar\_i}), and every $\varphi$-derivative (\code{phiN\_i})
to a concrete number. It is the \emph{universal substitution dictionary} ---
every downstream phase substitutes it into its symbolic expressions to get
numbers. If any symbol is left out of \code{num\_params}, the corresponding
numerical evaluation silently collapses (see Section~\ref{sec:co-gotchas}).
\end{defn}

\subsection{What \code{fsolve} actually does}

The numerical root-finding inside \code{solve\_mean\_field} uses SciPy. SciPy is
the standard scientific-computing library for Python; its \code{scipy.optimize}
submodule provides numerical root-finders and optimizers. The call
(\file{pipeline/\_mean\_field.py:174}) is:

\begin{lstlisting}
from scipy.optimize import fsolve              # _mean_field.py:11
attempt = fsolve(mf_residual, x0, full_output=True)
\end{lstlisting}

\code{fsolve} finds a root of the multivariate residual
$\texttt{mf\_residual}(n^{*}) = n^{*} - \text{RHS}(n^{*})$ --- the saddle
self-consistency written as ``input minus output''. \code{full\_output=True}
makes it return the tuple \code{(x, infodict, ier, msg)}, where \code{ier == 1}
signals clean convergence; the solver tries a sequence of initial guesses and
keeps the first that converges. This is the only place a Newton-type solver
appears in the temporal pipeline.

\section{Phase [3.5] --- the spatial-versus-temporal fork}\label{sec:co-spatial}

\textbf{Code:} \file{pipeline/compute.py:382}--\code{:626}.

Between the saddle and the pole-finder sits the single biggest branch in the
function. A spatial (PDE) model cannot go down the temporal road at all: its
propagator \code{G\_ft} carries an \emph{inert} \code{Laplacian} symbol that the
$\omega$-pole-finder and the time-domain Phase~J simply cannot consume. So
spatial models short-circuit here, routing to the momentum-space integrator and
returning early with a \emph{different} result dict.

\subsection{The branch logic}

First, two guard cases:

\begin{lstlisting}
if spatial_grid is not None and not model.get('spatial'):
    import warnings as _warnings
    _warnings.warn(
        'spatial_grid was provided but the model is not spatial ...', stacklevel=2)
\end{lstlisting}

Passing a \code{spatial\_grid} to a \emph{non}-spatial model is a no-op: the code
warns (\file{pipeline/compute.py:390}) and falls through to the temporal path.
The substantive branch is the converse --- the model \emph{is} spatial \emph{and}
a spatial input (\code{spatial\_grid} or \code{spatial\_points}) was given:

\begin{lstlisting}
if model.get('spatial') and (spatial_grid is not None
                             or spatial_points is not None):
    ...
\end{lstlisting}

Inside, the path splits on \code{k}.

\textbf{The $k=2$ grid path.} This is the workhorse spatial case: a two-point
correlator $C(x,\tau)$ on a grid of $x$-points. It requires \code{spatial\_grid}
(passing \code{spatial\_points} at $k=2$ raises a \code{ValueError} at
\file{pipeline/compute.py:414}, because \code{spatial\_points} is the $k\ge3$
event API). It guards \code{max\_ell $\le$ 2} and the stationary initial
condition, builds the $\tau$-grid and the $x$-grid, and then dispatches by loop
order: for \code{max\_ell $\ge$ 1} it calls
\code{compute\_spatial\_correlator\_generic} (the full-diagram momentum
integrator), and for tree level (\code{max\_ell == 0}) it calls
\code{compute\_spatial\_correlator\_via\_pipeline}. Crucially, the loop path
passes \code{parallel=(parallel and spatial\_parallel)}
(\file{pipeline/compute.py:511}) --- the two flags are \emph{and}-ed, and the
spatial parallelism is thread-based, not fork-based (Section~\ref{sec:co-parallel}).

It then builds a \code{total\_C} closure that evaluates the real-space correlator
at an arbitrary point, applying the 1-loop tadpole mass-shift correction
$\delta C = \Sigma\cdot\partial C_0/\partial A$ by finite difference for $d=1$
and by a radial inverse Fourier transform for $d\ge2$, and \textbf{returns early}
(\file{pipeline/compute.py:584}) a spatial dict whose shape we give in
Section~\ref{sec:co-result}.

\textbf{The $k\ge3$ event path.} For three or more legs the spatial output is not
a grid but a set of explicit \emph{events}, so it requires \code{spatial\_points}
--- an \code{(n\_pts, k-1, 2)} array of $(x_j, \tau_j)$ offsets. It calls
\code{compute\_spatial\_kpoint}; if that raises a \code{NotImplementedError}
mentioning \code{'single-field'}, the code falls through to the
\emph{coupled multi-field} driver, which builds reaction-diffusion matrices
$M, D, V$ from \code{prop['K\_ft']}, checks that the drift $V$ is zero
(\file{pipeline/compute.py:450}; nonzero drift raises), splits the reference
diffusion, and calls \code{compute\_coupled\_kpoint}. This path \textbf{returns
early} (\file{pipeline/compute.py:463}) a dict keyed by \code{C\_kpoint}.

\subsection{The clear-failure guard}

If a spatial model reaches the end of the branch without \emph{any} spatial input,
it would otherwise fall into the temporal pole-finder and die deep in Sage with a
cryptic \code{TypeError}. The code refuses to let that happen
(\file{pipeline/compute.py:619}):

\begin{lstlisting}
if model.get('spatial'):
    raise ValueError(
        f"model {model.get('name', '<unnamed>')!r} is spatial ...: "
        f"compute_cumulants requires spatial_grid=... to route to the "
        f"spatial integrator. ...")
\end{lstlisting}

\begin{gotcha}
The two spatial APIs are not interchangeable. \code{spatial\_grid} is the $k=2$
grid API and \code{spatial\_points} is the $k\ge3$ event API; mixing them raises.
A spatial $k\ge3$ run returns a different dict shape than a $k=2$ run (events, not
a $\tau$-grid), so any caller that handles both must branch on \code{k}. And a
spatial model with \emph{no} spatial input raises a clear \code{ValueError} rather
than a cryptic Sage error --- if you see that message, you forgot to pass
\code{spatial\_grid} (try the theory's \code{METADATA['spatial\_grid']}).
\end{gotcha}

The remainder of this chapter follows the \emph{temporal} path; the spatial
integrator has its own chapters --- this one is the junction, not the destination.

\section{Phase [4] --- numerical poles and residues}\label{sec:co-poles}

\textbf{Code:} \file{pipeline/compute.py:628}--\code{:633}.

With the saddle solved and \code{num\_params} in hand, the symbolic propagator
\code{G\_ft}$(\omega)$ can be made numeric and its poles extracted:

\begin{lstlisting}
compute_poles_and_residues(prop, num_params, verbose=verbose)
\end{lstlisting}

\code{compute\_poles\_and\_residues} (\file{pipeline/\_propagator.py:851}) uses a
three-tier strategy --- an exact polynomial fraction-field over
\code{CyclotomicField(4)} = $\mathbb{Q}[\ii]$, a numpy-cofactor fallback, and a
legacy symbolic path --- to find the retarded poles ($\operatorname{Im}\omega > 0$)
and the residue matrices. The exact path computes
$C_k[i,j] = \ii\, P_{ij}(\omega_k)/Q'(\omega_k)$.

\begin{gotcha}
\code{compute\_poles\_and\_residues} \emph{mutates \code{prop} in place}
(\file{pipeline/compute.py:632}). There is no return value used: the orchestrator
relies entirely on the side effect of filling \code{prop['pole\_vals']} and
\code{prop['C\_mats']} (and \code{D\_delta}, if phase [2] left it \code{None}).
The very same \code{prop} object is later sliced into \code{propagator\_data} for
phase [7] and returned in the result dict. If you are reading the code top to
bottom, do not expect an assignment here --- the data appears inside \code{prop}.
\end{gotcha}

\section{Phase [5] --- enumerating the diagrams}\label{sec:co-diagrams}

\textbf{Code:} \file{pipeline/compute.py:635}--\code{:664}.

This phase produces the list of Feynman diagrams to integrate. First it builds
the field-index maps that tell the enumerator which ring variable is which field:

\begin{lstlisting}
ring_var_names = list(ft._ns._ring_var_names)
n_tilde = ft._n_tilde
resp_idx, phys_idx = build_field_index_map(ring_var_names, n_tilde)
for f in external_fields:
    if f not in phys_idx:
        raise ValueError(
            f'external field {f} not in phys_idx ...')
\end{lstlisting}

\code{build\_field\_index\_map} splits the ring's generator names into the
response block (the first \code{n\_tilde}) and the physical block, returning two
dicts \code{resp\_idx} and \code{phys\_idx}. The validation loop catches a mistyped
external field early --- before enumeration --- with a helpful message listing the
valid keys. Then the enumerator runs:

\begin{lstlisting}
unique_by_ell, multiplicity_by_ell, all_unique = enumerate_unique_diagrams(
    ft, model,
    k=k, max_ell=max_ell,
    external_fields=external_fields,
    G_ft=prop['G_ft'],
    resp_idx=resp_idx, phys_idx=phys_idx,
    vtypes=vtypes, stypes=stypes,
    use_cache=use_cache, parallel=parallel,
    n_workers=n_workers, verbose=verbose,
)
\end{lstlisting}

\code{enumerate\_unique\_diagrams} (\file{pipeline/\_diagrams.py:54}) returns the
dicts \code{unique\_by\_ell} ($\{\ell:$ list of \code{TypedDiagram}$\}$) and
\code{multiplicity\_by\_ell} ($\{\ell:$ list of int$\}$) plus the flat
concatenation \code{all\_unique}. For each loop order $\ell$ it checks a disk
cache and, on a miss, runs the four-stage diagram pipeline:

\begin{defn}
\textbf{Prediagram $\to$ typed $\to$ causal $\to$ unique.}
The four-stage enumeration is: \term{prediagram} enumeration (topological
skeletons), \code{enumerate\_all\_typed} (assigning a definite field type to every
leg --- this is the parallelizable stage, fork-based when \code{parallel=True}),
\code{filter\_causal} (dropping acausal diagrams), and
\code{deduplicate\_with\_multiplicities} (collapsing isomorphic graphs and
recording how many collapsed into each class).
\end{defn}

\begin{gotcha}
The cache stage name embeds a version suffix: \code{unique\_typed\_mult\_v3\_...}
(\file{pipeline/\_diagrams.py:150}). The \code{v3} is not decoration. Three
documented bumps lie behind it --- multiplicity-aware dedup, the
\code{NoiseSourceType} promotion, and the \emph{complete}-isomorphism signature.
Loading an older cache would silently resurrect a real bug (non-isomorphic
diagram classes collided and had their integrals dropped at $k\ge3$), so the
version suffix forces a rebuild rather than trusting stale data. The expand cache
and the propagator cache live separately, under \file{saved\_theories/<tag>/}.
\end{gotcha}

\section{Phase [6] --- symmetry factors and prefactors}\label{sec:co-classify}

\textbf{Code:} \file{pipeline/compute.py:666}--\code{:705}.

Each diagram needs a scalar weight before it can be integrated. This phase walks
the diagrams \emph{by loop order} (so each record carries its $\ell$ tag, which
phase [7] needs) and classifies each one:

\begin{lstlisting}
for ell in sorted(unique_by_ell.keys()):
    mults = multiplicity_by_ell.get(ell, [1] * len(unique_by_ell[ell]))
    for td, mult in zip(unique_by_ell[ell], mults):
        info = classify_coefficient_factors(
            td, time_dep_params, noise_structure)
        combined_prefactor = SR(info['scalar_prefactor'])
        kernel_groups.append({
            'diagrams':           [td],
            'combined_prefactor': combined_prefactor,
        })
        diagram_records.append({
            'typed_diagram':      td,
            'classify':           info,
            'combined_prefactor': combined_prefactor,
            'multiplicity':       mult,
            'ell':                ell,
        })
\end{lstlisting}

\code{classify\_coefficient\_factors} (\file{msrjd/diagrams/symmetry.py:618})
returns a dict with the keys \code{Scal}, \code{scalar\_prefactor},
\code{vertex\_time\_factors}, \code{source\_time\_info}, and \code{is\_stationary}.
The orchestrator pulls \code{scalar\_prefactor} out, wraps it with \code{SR(...)}
to coerce it to a symbolic expression, and appends to two parallel lists:
\code{kernel\_groups} (the legacy Phase-J input format) and \code{diagram\_records}
(the richer per-diagram record carried in the result).

\subsection{The symmetry factor \texorpdfstring{$\Sym(\Gamma)$}{S(Gamma)}}

The mathematical content of this phase is the combinatorial weight of each
diagram. A connected $k$-point cumulant is a weighted sum of diagrams; each
diagram $\Gamma$ contributes
\begin{equation}
  \mathrm{weight}(\Gamma)
  \;=\;
  \Sym(\Gamma)\,
  \prod_v c_{\mathrm{user},v}\,
  \prod_e P_e\,
  \int \prod_v \dd t_v,
  \label{eq:co-weight}
\end{equation}
where $c_{\mathrm{user},v}$ is the literal coefficient written in the action (no
implicit $1/n!$), $P_e$ is the propagator on edge $e$, the integral runs over the
internal vertex times, and $\Sym(\Gamma)$ is the combinatorial symmetry factor.
\Daedalus{} uses the canonical Feynman rule (``Path~A'' in this codebase):
\begin{equation}
  \Sym(\Gamma)
  \;=\;
  \frac{\displaystyle\prod_v \prod_\ell n_{v,\ell}!}
       {\bigl|\Aut_{\text{fixed-ext}}(\Gamma)\bigr|}.
  \label{eq:co-Sym}
\end{equation}
The numerator is the per-vertex Wick combinatorial --- factorials of
identical-leg multiplicities at each vertex, for both response and physical legs.
The denominator is the order of the automorphism group of the colour-preserving
incidence digraph, with the external leaves \emph{pinned} at their canonical
positions. This single ratio accounts for every Feynman-rule symmetry at once:
same-type vertex swaps, parallel-edge swaps, self-loop leg swaps.

The denominator is computed by Sage. \code{combinatorial\_factor}
(\file{msrjd/diagrams/symmetry.py:403}) computes the Wick numerator via
\code{\_wick\_leg\_factor} (\file{msrjd/diagrams/symmetry.py:115}) and the
automorphism order via \code{\_automorphism\_order}
(\file{msrjd/diagrams/symmetry.py:328}), which in turn calls Sage's
\code{D.automorphism\_group(partition=...)}
(\file{msrjd/diagrams/symmetry.py:339}). Sage routes that call through its bundled
graph-canonicalization backend (the \code{bliss} / \code{nauty}-family
algorithms), and \code{.order()} of the returned permutation group is
$|\Aut_{\text{fixed-ext}}(\Gamma)|$.

\begin{note}
\code{nauty} (``No AUTomorphisms, Yes?'') and \code{bliss} are C libraries that
compute graph automorphism groups and canonical forms --- the gold-standard tools
for graph isomorphism. You give them a graph with an optional vertex colouring and
they return either the automorphism group or a canonical labeling such that two
graphs are isomorphic if and only if their canonical forms are byte-identical.
\Daedalus{} never calls them directly; Sage's \code{DiGraph.automorphism\_group}
and \code{DiGraph.canonical\_label} (\file{msrjd/diagrams/symmetry.py:517}) wrap
them. The same machinery powers the dedup signature in phase [5]: the
\emph{canonical form} of the coloured incidence digraph is a complete isomorphism
invariant.
\end{note}

\begin{gotcha}
The dedup \code{multiplicity} from phase [5] is recorded in each record but
\emph{not multiplied} into the prefactor (\file{pipeline/compute.py:686}--\code{:692}).
Under Path~A the symmetry factor $\Sym(\Gamma)$ in \eqref{eq:co-Sym} already
carries each equivalence class's full weight via the orbit--stabilizer count, so
multiplying by the class size would double-count. \code{multiplicity} is kept
purely for diagnostics. (Historically a caller-side multiplication compensated for
an \emph{incomplete} signature that merged non-isomorphic diagrams; now that the
signature is a complete invariant, that compensation is obsolete and removed.)
\end{gotcha}

\subsection{What \code{classify\_coefficient\_factors} does internally}

For completeness, the per-diagram classifier (\file{msrjd/diagrams/symmetry.py:618})
does the following for one \code{TypedDiagram}:

\begin{enumerate}
  \item Computes \code{Scal = combinatorial\_factor(td)} --- the symmetry factor
        $\Sym(\Gamma)$.
  \item Walks every vertex; each vertex coefficient acquires a $(-1)$ because the
        weight is $\int \exp(-S)$ (\code{coeff = -SR(vtype.coefficient)},
        \file{msrjd/diagrams/symmetry.py:643}).
  \item For \emph{source (noise) vertices}, decides whether the amplitude can be
        pulled out of the time integral --- a time-independent white amplitude
        pulls out into the scalar prefactor; otherwise it stays in
        \code{source\_time\_info} for Phase~J.
  \item For \emph{interaction vertices}, splits the coefficient into a
        time-independent constant part (joining \code{scalar\_parts}) and a
        time-dependent part (kept in \code{vertex\_time\_factors}).
  \item Multiplies all the constant parts into \code{scalar\_prefactor}.
  \item Sets \code{is\_stationary} \code{True} iff there are no per-vertex time
        factors and no time-dependent noise amplitudes --- which tells Phase~J it
        may use the fast time-translation-invariant path.
\end{enumerate}

\section{Phase [7] --- Phase J and the per-\texorpdfstring{$\ell$}{ell} assembly}\label{sec:co-phasej}

\textbf{Code:} \file{pipeline/compute.py:707}--\code{:799}.

The final phase performs the actual loop integrals. It first builds the
$\tau$-grid and the \code{propagator\_data} dict --- precisely the subset of
\code{prop} that Phase~J consumes (\file{pipeline/compute.py:714}):

\begin{lstlisting}
tau_grid = (np.asarray(tau_grid, dtype=float) if tau_grid is not None
            else np.arange(-tau_max, tau_max + tau_step * 0.5, tau_step))

propagator_data = {
    'K_ker':   prop['K_ker'],   'K_ft':    prop['K_ft'],
    'G_ft':    prop['G_ft'],    'adj_ft':  prop['adj_ft'],
    'D_omega': prop['D_omega'], 'D_delta': prop['D_delta'],
    't_var':   prop['t_var'],   'omega':   prop['omega'],
    'nf':      prop['nf'],
    'pole_vals': prop['pole_vals'],  'C_mats': prop['C_mats'],
}
\end{lstlisting}

Note that \code{tau\_grid} runs from $-\texttt{tau\_max}$ to $+\texttt{tau\_max}$;
the \code{tau\_step * 0.5} nudge in \code{np.arange}'s upper bound is the standard
trick to make the endpoint inclusive despite floating-point round-off.

\subsection{Choosing the \texorpdfstring{$\tau$}{tau}-evaluation pattern}

The grid is evaluated differently depending on the number of external legs
(\file{pipeline/compute.py:729}):

\begin{lstlisting}
if k == 1:
    tau_points = [(float(t),) for t in tau_grid]
elif k == 2:
    tau_points = [(0.0, float(t)) for t in tau_grid]  # leaf 0 pinned, leaf 1 swept
else:
    tau_points = None   # k>=3: caller handles grid evaluation
\end{lstlisting}

For a one-point function each grid point is a single time; for a two-point
function leaf $0$ is pinned to $0$ and leaf $1$ is swept over the grid (the
familiar single-axis correlation slice $C(\tau)$); for $k\ge3$ no automatic grid
eval is done --- \code{tau\_points = None} and the caller supplies the evaluation
points.

\subsection{Once per loop order}

The heart of the phase is a loop over $\ell$. Running Phase~J \emph{separately for
each loop order} is what gives the per-order decomposition the notebooks plot as
``tree'', ``1-loop'', ``2-loop'':

\begin{lstlisting}
for ell in sorted(unique_by_ell.keys()):
    records_ell = [r for r in diagram_records if r['ell'] == ell]
    if not records_ell:
        total_C_by_ell[ell] = (lambda *t: 0.0 + 0.0j)
        phase_j_by_ell[ell] = None
        C_tau_by_ell[ell] = (np.zeros(len(tau_grid), dtype=complex)
                             if tau_points is not None else None)
        continue
    ...
    td_result_ell = compute_correction_td(
        typed_diagrams = [r['typed_diagram'] for r in records_ell],
        prefactors     = [r['combined_prefactor'] for r in records_ell],
        k=k, propagator_data=propagator_data,
        external_fields=external_fields, num_params=num_params,
        origin_leaf_idx=origin_leaf_idx,
    )
    total_C_by_ell[ell] = td_result_ell['total_C']
    phase_j_by_ell[ell] = td_result_ell
    if tau_points is not None:
        C_tau_by_ell[ell] = np.array(
            td_result_ell['total_C_batch'](
                tau_points, parallel=parallel, n_workers=n_workers),
            dtype=complex)
\end{lstlisting}

\code{compute\_correction\_td} (\file{msrjd/integration/time\_domain/pipeline.py:202})
is the per-$\ell$ workhorse: for each typed diagram it numerically integrates the
vertex times against the pole-residue propagator and accumulates a contribution
callable. It returns a dict with \code{total\_C} (a callable summing the diagrams
serially) and \code{total\_C\_batch} (a callable that fans the $\tau$-grid out
over a fork-based process pool, or degrades to serial under the notebook fork
guard --- Section~\ref{sec:co-parallel}). The orchestrator stores the per-$\ell$
callable, the full per-$\ell$ result dict, and --- for $k\in\{1,2\}$ --- the
batch-evaluated array over the grid.

\begin{gotcha}
When a loop order has \emph{no} diagrams, the code still inserts a zero callable
and a zero array (\file{pipeline/compute.py:746}). This is necessary, not
defensive padding: some $(k,\ell)$ configurations genuinely have no diagrams at a
given order, and without the placeholder the master \code{total\_C} sum below
would be missing a key.
\end{gotcha}

\begin{note}
Setting \code{use\_grouped\_phase\_j=True} routes this call to
\code{compute\_correction\_td\_grouped} (\file{pipeline/compute.py:761}) instead
of \code{compute\_correction\_td}. The grouped integrator is a prototype that
sums integrands grouped by parent prediagram before quadrature; see
\file{pipeline/\_grouped\_phase\_j.py} for its math and caveats. The default path
is the ungrouped \code{compute\_correction\_td}.
\end{note}

\subsection{Summing across loop orders}

After the per-$\ell$ loop, the master correlator is assembled by summation
(\file{pipeline/compute.py:802}):

\begin{lstlisting}
def total_C(*ext_time_values):
    return sum(complex(fn(*ext_time_values))
               for fn in total_C_by_ell.values())

if tau_points is not None:
    C_tau = sum(C_tau_by_ell[ell] for ell in C_tau_by_ell)
else:
    C_tau = None
\end{lstlisting}

The master callable \code{total\_C} is the additive sum of the per-$\ell$
callables, and the master grid array \code{C\_tau} is the elementwise sum of the
per-$\ell$ grid arrays:
\begin{equation}
  \texttt{total\_C}(\tau) \;=\; \sum_\ell C_\ell(\tau),
  \qquad
  \texttt{C\_tau} \;=\; \sum_\ell \texttt{C\_tau\_by\_ell}[\ell].
  \label{eq:co-sum}
\end{equation}

\begin{note}
\textbf{On the phrase ``set-partition.''}
You may see the framing ``moment / cumulant / central-moment set-partition
assembly'' elsewhere. The orchestrator itself does \emph{not} contain any
moment$\leftrightarrow$cumulant set-partition (M\"obius-inversion) bookkeeping.
\code{compute\_cumulants} returns the \emph{connected} cumulant directly --- the
diagrammatic sum \emph{is} the cumulant. The only assembly this file performs is
the additive per-$\ell$ sum in \eqref{eq:co-sum}. Moment and central-moment
conversions, when needed, live in downstream consumers (the engine-API and
cumulants/moments chapters), not in \file{pipeline/compute.py}.
\end{note}

\section{The adaptive mean-field dictionary}\label{sec:co-mfdict}

\textbf{Code:} \file{pipeline/compute.py:812}--\code{:879}.

Before packing the result, the function discovers which model parameters are
mean-field saddles and reads their solved values into a tidy
\code{\{internal\_name: [v\_pop1, ...]\}} dict. It tries three sources in order:
the \code{naming\_convention['mf\_parameters']} list (preferred, written by the
\code{TheoryBuilder} when \code{mean\_field=True} is set on a parameter), the
per-parameter \code{mean\_field} flag, and finally the classic
\code{nstar}/\code{vstar}/\code{mstar} fallback for hand-written models that
pre-date the declarative API.

For each discovered saddle name it maps the name to the range of population
indices it owns --- via the closure \code{\_saddle\_indices}, which handles
heterogeneous theories that declare \code{indexed\_by=['<pop>']} on each saddle ---
and reads each value out of \code{num\_params}.

\begin{gotcha}
Entries that are \emph{entirely} NaN are dropped (\file{pipeline/compute.py:878}).
A model may declare a saddle symbol it never actually substitutes --- e.g.
\code{mstar} in a non-GTaS model that nonetheless ships an \code{mstar}
declaration. Without this filter those would surface in the result as all-NaN
vectors. The check \code{not all(v != v for v in vals)} is the idiomatic
``not every entry is NaN'' test (recall \code{nan != nan} is \code{True}).
\end{gotcha}

\section{The result dict}\label{sec:co-result}

\textbf{Code:} \file{pipeline/compute.py:881}--\code{:908}.

Everything the run produced is packed into one dictionary. For the temporal path
the keys are:

\begin{center}
\small
\begin{longtable}{@{}p{0.27\textwidth}p{0.20\textwidth}p{0.45\textwidth}@{}}
\toprule
\textbf{key} & \textbf{type} & \textbf{meaning} \\
\midrule
\endhead
\code{total\_C} & callable & $\sum_\ell C_\ell(*\tau)$ \\
\code{total\_C\_by\_ell} & $\{\ell:$ callable$\}$ & per-loop-order callable \\
\code{C\_tau} & ndarray or \code{None} & total on the $\tau$-grid ($k\in\{1,2\}$; \code{None} for $k\ge3$) \\
\code{C\_tau\_by\_ell} & $\{\ell:$ ndarray$\}$ & per-loop-order grid; \code{C\_tau} is their sum \\
\code{tau\_grid} & ndarray & the $\tau$ values \\
\code{mf\_values} & $\{$name: $[v_1,\dots]\}$ & adaptive saddle values \\
\code{mf} & \code{MeanField} & natural-name saddle accessor \\
\code{params} & \code{Parameters} & natural-name parameter accessor \\
\code{num\_params} & $\{$\code{SR}: float$\}$ & numerical substitution dict \\
\code{propagator} & dict & the \code{prop} dict (poles/residues filled) \\
\code{diagrams} & list & the \code{diagram\_records} \\
\code{kernel\_groups} & list & per-diagram Phase-J input format \\
\code{phase\_j\_by\_ell} & $\{\ell:$ dict$\}$ & per-$\ell$ Phase-J result dicts \\
\code{phase\_walls} & $\{$label: sec$\}$ or \code{None} & per-phase wall-time \\
\code{config} & dict & echoed args (internal + user external fields, taylor order, \dots) \\
\bottomrule
\end{longtable}
\end{center}

The two natural-name accessors deserve a sentence. \code{mf} is a
\code{MeanField} object (\file{pipeline/access.py}): \code{mf['v', 1]} returns
$v^{*}_1$ and \code{mf['n']} returns the whole \code{nstar} vector.
\code{params} is a \code{Parameters} object: \code{params['E', 1]} and
\code{params['w', 1, 2]} index into vector and matrix parameters with 1-based
indexing. These spare downstream code from ever touching the raw symbolic dict.

\subsection{Two early-return shapes}\label{sec:co-datastructures}

The spatial paths each return a \emph{different} dict, earlier in the function.
The spatial $k=2$ early return (\file{pipeline/compute.py:584}) adds
\code{C\_tau\_x} (the full $(n_\tau, n_x)$ real-space correlator),
\code{C\_tau\_x\_by\_order} (cumulative per-loop-order grids),
\code{spatial\_grid}, and \code{spatial\_info}; its \code{total\_C} is the spatial
closure (1 arg $\to$ $C(\tau)$ at $x=0$, 2 args $\to$ $C(x,\tau)$) and its
\code{total\_C\_by\_ell} is \code{\{0: total\_C\}}. The spatial $k\ge3$ early
return (\file{pipeline/compute.py:463}) is keyed instead by \code{C\_kpoint} (an
\code{(n\_pts,)} array), \code{C\_kpoint\_by\_ell}, \code{points} (the echoed
\code{spatial\_points}), \code{spatial\_info}, \code{k}, \code{max\_ell}, and
\code{mf}.

For reference, the \code{prop} dict that phases [2] and [4] build and fill has the
following keys (the propagator chapter is the authority):

\begin{center}
\small
\begin{longtable}{@{}p{0.18\textwidth}p{0.30\textwidth}p{0.44\textwidth}@{}}
\toprule
\textbf{key} & \textbf{type} & \textbf{meaning} \\
\midrule
\endhead
\code{K\_ker}   & Sage \code{matrix(SR)} & bilinear kernel in time ($\delta$, $\delta'$) \\
\code{K\_ft}    & Sage \code{matrix(SR)} & Fourier image of \code{K\_ker} (rational in $\omega$) \\
\code{G\_ft}    & Sage matrix or \code{None} & propagator \code{K\_ft$^{-1}$}; \code{None} for rich matrices \\
\code{adj\_ft}  & Sage matrix or \code{None} & adjugate $=$ \code{G\_ft}$\cdot$det \\
\code{D\_omega} & SR or \code{None} & $\det(\text{\code{K\_ft}})$ \\
\code{D\_delta} & Sage \code{matrix(SR)} & instantaneous part $\lim_{\omega\to\infty}$\code{G\_ft} \\
\code{t\_var}, \code{omega} & SR vars & the time and frequency symbols \\
\code{nf}       & int & number of field components \\
\code{pole\_vals} & list[complex] & retarded poles, filled by phase [4] \\
\code{C\_mats}  & list[Sage \code{matrix(CDF)}] & residue matrices, filled by phase [4] \\
\bottomrule
\end{longtable}
\end{center}

\subsection{DAE extras, and the save}

\textbf{Code:} \file{pipeline/compute.py:910}--\code{:944}. When the DAE
equation-based solver ran in phase [3], the result picks up multi-root extras:
\code{mf\_all\_roots} (every distinct root with a stability annotation),
\code{mf\_stable\_roots}, \code{mf\_unstable\_roots}, \code{mf\_index\_used},
\code{mf\_state\_var\_order}, and \code{mf\_stability}.

\begin{gotcha}
\code{fixed\_point\_index} indexes the \emph{stable} subset only
(\file{pipeline/compute.py:910} comment). \code{mf\_stable\_roots} is that subset
in sort order; unstable roots are surfaced via \code{mf\_unstable\_roots} for
inspection but are \emph{not} selectable as expansion points.
\end{gotcha}

Finally, if \code{output\_npz} or \code{output\_csv} was given, the function calls
\code{save\_npz} / \code{save\_csv} from \file{pipeline/save.py} (that module owns
the on-disk schema --- per-loop-order curves \code{C\_tree}, \code{C\_1\_loop},
\dots, the adaptive mean-field arrays, the parameter keys).

\section{Parallelism: fork versus thread, and the macOS hazard}\label{sec:co-parallel}

\Daedalus{} has two parallelizable stages on the temporal path --- per-prediagram
type assignment in phase [5] and per-$\tau$ Phase~J evaluation in phase [7] ---
and both are governed by the single \code{parallel} flag. There is also a
separate \code{spatial\_parallel} flag for the spatial loop integrator. The
distinction between them is not cosmetic; it is a hard-won safety boundary.

\begin{defn}
\textbf{Fork versus thread parallelism.}
\term{Fork} parallelism copies the entire parent process (it is fast to start,
but on macOS it crashes a Jupyter kernel that has already initialised Cocoa, BLAS,
or matplotlib). \term{Thread} parallelism shares one process's memory (it is safe,
and numpy releases the GIL during its heavy linear algebra so threads genuinely
run in parallel). The temporal stages \emph{fork}; the spatial stages
\emph{thread}.
\end{defn}

\begin{gotcha}
\textbf{Spatial parallelism is thread-based, never fork-based.} The default
\code{spatial\_parallel=True} routes the loop integrator through a
\code{ThreadPoolExecutor}. Forking after Cocoa/BLAS/matplotlib initialisation
inside a macOS Jupyter kernel \emph{hard-crashes the kernel and the OS}, and
setting \code{OBJC\_DISABLE\_INITIALIZE\_FORK\_SAFETY=YES} makes it \emph{worse},
not better. The temporal \code{parallel} flag \emph{is} fork-based, but it is
guarded: \code{fork\_unsafe\_in\_notebook} (\file{msrjd/fork\_safety.py:27}) trips
only on \code{darwin} \emph{and} a live ZMQ/Jupyter kernel \emph{and} the
\code{fork} start method, degrading that single case to serial with a one-time
warning. Under pytest, on Linux, in a terminal, or from a script, fork is kept.
This is why the example notebooks in Chapter~\ref{ch:quickstart} pass
\code{parallel=False} explicitly: it is the safe default in a notebook, and the
tiny example models are fast enough serially.
\end{gotcha}

\section{Performance: the phase-wall summary}\label{sec:co-perf}

\textbf{Code:} \file{pipeline/compute.py:946}--\code{:955}. When verbose, the
function closes with a one-line \code{Done.} and, if it tracked timings, a
per-phase breakdown:

\begin{lstlisting}
if verbose:
    print(f'\nDone.  k={k}, max_ell={max_ell}, '
          f'{len(all_unique)} unique diagrams, {len(tau_grid)} tau points.')
    if phase_walls:
        total = sum(phase_walls.values())
        print(f'\n  Phase wall summary (Sigma = {total:.1f}s):')
        for label, dt in phase_walls.items():
            pct = 100.0 * dt / total if total > 0 else 0.0
            print(f'    {label:12s}  {dt:6.2f}s  ({pct:4.1f}%)')
\end{lstlisting}

This reads back the \code{phase\_walls} dict the \code{\_phase\_time} closure has
been filling all along, prints each phase's seconds and its percentage of the
total, and returns the \code{result}. In nearly every nontrivial run the phase
that dominates the clock is \code{phase\_j} (phase [7]) --- the loop integration
--- which is why so much of this manual is devoted to it. The summary is your
first diagnostic when a run is slow: it tells you immediately whether you are
paying for the symbolic expand, the pole-finder, the diagram enumeration, or the
integration.

\begin{gotcha}
The phase-wall summary is opt-in: \code{phase\_walls} is \code{None} when
\code{verbose=False} (\file{pipeline/compute.py:296}), so the timing dict in the
result is only populated in a verbose run. A quiet run returns
\code{result['phase\_walls'] = None}.
\end{gotcha}

\section{The data flow, threaded between phases}\label{sec:co-dataflow}

It is worth seeing the whole thread of data in one diagram, because the
orchestrator's real subject is \emph{what each phase hands the next}:

\begin{lstlisting}[style=console]
  model --+--------------------------------------------------------------+
          v                                                              |
  [1] FieldTheory(model, taylor_order).expand() -> ft (_ns, _n_tilde,    |
          |                                            sectors, vtypes)   |
          v                                                              |
  [2] build_propagator(ft, model) -> prop {K_ft, G_ft, D_delta, w, nf}   |
          |                                                              |
          v                                                              |
  [3] solve_mean_field(ft, model, fundamental) -> num_params {SR:float}  |
          |                                                              |
          v  (spatial? -> early return via the momentum integrator)      |
  [4] compute_poles_and_residues(prop, num_params) -> prop.pole_vals,    |
          |                                              prop.C_mats      |
          v                                                              |
  [5] enumerate_unique_diagrams(ft, model, k, ell, ext, G_ft,           |
          |   resp/phys_idx, vtypes, stypes) -> unique_by_ell            |
          v                                                              |
  [6] classify_coefficient_factors(td, time_dep, noise) ->              |
          |   diagram_records [{td, classify, prefactor, mult, ell}]     |
          v                                                              |
  [7] compute_correction_td(td_list[ell], prefactors[ell],             |
          |   propagator_data, k, external_fields, num_params)           |
          |   -> td_result {total_C, total_C_batch, ...}                 |
          v                                                              |
      total_C = sum_ell ;  C_tau = sum_ell  ->  result dict -------------+
\end{lstlisting}

\begin{example}[A Hawkes-like $k=2$ tree run]
Take \code{external\_fields = [('n',1),('n',2)]}. At validation
(\file{pipeline/compute.py:288}) it is normalized to \code{[('dn',1),('dn',2)]}.
With \code{max\_ell = 0} the Taylor budget \eqref{eq:co-taylor} is
$\max(2,2)=2$. Phases [1]--[6] run; phase [5] finds the tree-level diagrams; phase
[7] sets \code{tau\_points = [(0.0, t) for t in tau\_grid]} (leaf 0 pinned to 0,
leaf 1 swept). Each $\tau$-point is passed as \code{total\_C(0.0, t)}, producing
\code{C\_tau\_by\_ell[0]} as a complex array. With no loops, \code{C\_tau} equals
\code{C\_tau\_by\_ell[0]} exactly. The notebook then plots \code{tau\_grid}
against \code{C\_tau.real} --- the autocorrelation you saw in
Chapter~\ref{ch:quickstart}.
\end{example}

\section{The two helper functions}\label{sec:co-helpers}

The file defines two small helpers above the entry point, both pure formatting
with no effect on the computation.

\code{\_trunc(s, maxlen=200)} (\file{pipeline/compute.py:45}) returns
\code{str(s)} truncated to \code{maxlen} characters with a trailing \code{'...'}.
It keeps long symbolic coefficients readable in the verbose printout.

\code{\_print\_action\_sectors(ft)} (\file{pipeline/compute.py:50}) pretty-prints
the three action sectors after expansion: the mean-field sectors $(0,0)$,
$(1,0)$, $(0,1)$ (read from \code{ft.\_mf\_sector\_raw}), the free $(1,1)$
bilinear kernel (\code{ft.free\_action()}), and every interaction sector of total
degree $\ge2$ excluding $(1,1)$ (iterating \code{ft.sectors()}). Its inner helper
\code{\_fmt\_poly} walks a Sage polynomial's \code{.dict()} --- the
exponent-vector $\to$ coefficient map --- and reconstructs each monomial string
from the ring generator names. This is the ``here is your action'' diagnostic; it
prints, and nothing else.

\section{Gotchas and failure modes}\label{sec:co-gotchas}

We have flagged the local gotchas inline; here is the consolidated list of the
ways a \code{compute\_cumulants} run can surprise you.

\begin{itemize}
  \item \textbf{\code{num\_params} must be fully numeric.} If the mean-field phase
        leaves \emph{any} kernel symbol or saddle value symbolic, the pole-finder
        silently returns zero roots and the whole correlator collapses to $0$ ---
        a subtle bug where, for instance, a missing \code{phi0\_i} (resolved from
        \code{mf\_bg\_conditions}, \file{pipeline/\_mean\_field.py:272}) zeroes the
        result. The pole-finder warns when it finds no roots
        (\file{pipeline/\_propagator.py:976}), so a sudden flat-zero correlator is
        the symptom to watch.
  \item \textbf{\code{prop} is mutated in place by phase [4].} The poles and
        residues appear \emph{inside} the \code{prop} object, not as a return
        value (Section~\ref{sec:co-poles}).
  \item \textbf{Spatial models must be given the right spatial input.}
        \code{spatial\_grid} for $k=2$, \code{spatial\_points} for $k\ge3$; the
        wrong one (or neither) raises a clear error rather than a cryptic Sage one
        (Section~\ref{sec:co-spatial}).
  \item \textbf{Stale caches.} The diagram cache version suffix
        (\code{...\_v3\_...}) forces a rebuild after the dedup-signature fixes;
        the expand and propagator caches are separate
        (Section~\ref{sec:co-diagrams}).
  \item \textbf{The Taylor floor of $2$.} Probing higher-order vertices for
        cache-invalidation testing needs an explicit \code{taylor\_order} override
        (Section~\ref{sec:co-args}).
  \item \textbf{The grouped Phase~J is a prototype.}
        \code{use\_grouped\_phase\_j=True} routes to an experimental integrator;
        the default and trusted path is \code{compute\_correction\_td}
        (Section~\ref{sec:co-phasej}).
\end{itemize}

\section{What you now understand}\label{sec:co-summary}

\code{compute\_cumulants} is the one function that turns a model declaration plus
a parameter point into a numerical correlation function, and it does so by
sequencing seven phases: expand the action, build the propagator, solve the
saddle, find the poles, enumerate the diagrams, classify their prefactors, and
integrate them in the time domain --- with a hard fork between the temporal and
spatial routes after the saddle. It performs almost no arithmetic itself; its
real content is the \emph{contracts} between phases and the \emph{glue} that
threads \code{ft}, \code{prop}, \code{num\_params}, and \code{diagram\_records}
from one to the next, the additive per-$\ell$ assembly that produces the
loop-resolved correlator, and the disciplined parallelism that keeps a macOS
notebook from crashing. Every phase named here has its own chapter; this one is
the map that shows how they connect. The next chapters move outward from the
engine to the user surface --- the \code{daedalus} front desk that you actually
type at --- and then to the cumulant/moment conversions and the simulators that
validate everything the engine produces.
